\documentclass[11pt,a4paper]{article}
\usepackage[margin=21mm,headheight=16pt]{geometry}
\usepackage{fontspec}\setmainfont{Noto Serif}\setsansfont{Noto Sans}\setmonofont{Noto Sans Mono}
\usepackage{polyglossia}\setdefaultlanguage{turkish}
\usepackage{amsmath,amssymb,mathtools,graphicx,xcolor,booktabs,tabularx,array,enumitem}
\usepackage[most]{tcolorbox}\usepackage{fancyhdr,titlesec,needspace}
\usepackage[unicode,hidelinks,breaklinks]{hyperref}\usepackage{bookmark}
\definecolor{ink}{HTML}{18233A}\definecolor{accent}{HTML}{545AC4}\definecolor{green}{HTML}{087E66}\definecolor{paper}{HTML}{F3F5FA}
\color{ink}\setlength{\parindent}{0pt}\setlength{\parskip}{7pt}\linespread{1.10}
\setlength{\emergencystretch}{3em}\allowdisplaybreaks
\setlist{topsep=4pt,itemsep=4pt,leftmargin=*}
\pagestyle{fancy}\fancyhf{}\fancyhead[L]{\sffamily\small MÜHENDİSLİK MATEMATİĞİ}\fancyhead[R]{\sffamily\small Ders çalışma notları}\fancyfoot[L]{\sffamily\scriptsize Fotoğraflar + müfredat | 8 Kasım 2026}\fancyfoot[R]{\sffamily\thepage}
\titleformat{\section}{\Large\sffamily\bfseries\color{accent}}{\thesection}{.6em}{}
\titleformat{\subsection}{\large\sffamily\bfseries}{\thesubsection}{.6em}{}
\tcbset{enhanced,breakable,boxrule=.4pt,arc=2mm,left=3mm,right=3mm,top=2mm,bottom=2mm,before skip=8pt,after skip=8pt,fonttitle=\sffamily\bfseries,colback=paper,colframe=accent,coltitle=white}
\newcommand{\notebox}[2]{\begin{tcolorbox}[title={#1},colframe=accent!70!black]#2\end{tcolorbox}}
\hypersetup{pdfauthor={Mühendislik Matematiği Çalışma Atölyesi},pdfsubject={Ders fotoğrafları ve altı başlıklı müfredata göre özgün Türkçe çalışma notları}}
\begin{document}

\hypersetup{pdftitle={Soru çözüm atölyesi}}
{\sffamily\small\color{accent} ÇALIŞMA ATÖLYESİ / ÇÖZÜMLÜ SORULAR\par}
\vspace{8mm}
{\Huge\sffamily\bfseries\raggedright\hyphenpenalty=10000 Soru çözüm atölyesi\par}
\vspace{3mm}
{\large Ders fotoğraflarından alınan örnekler ve ek pekiştirme soruları. Her soruda tam çözüm ve kaynak etiketi.\par}
\vspace{4mm}\hrule\vspace{4mm}89 çözümlü örnek. Çıkmış sınav sorusu içermez. Birden fazla alt sorusu bulunan örnekler tek kayıt olarak sayıldı.
\tableofcontents\clearpage
\clearpage\section{Üç boyutlu uzay ve vektörler}
\Needspace{10\baselineskip}\begin{tcolorbox}[title={Örnek 1.1 | Birim vektör},colframe=accent]\(\mathbf u=(3,4)\) yönündeki birim vektörü bulunuz.
\par\vspace{1mm}{\scriptsize\sffamily Ders fotoğrafından | \nolinkurl{IMG-20261009-WA0026.jpg}}\end{tcolorbox}
\begin{tcolorbox}[title={Adım adım çözüm},colback=green!3!white,colframe=green]\begin{enumerate}
\item Önce norm alınır: \(\|\mathbf u\|=\sqrt{9+16}=5\).
\item Her bileşeni 5 ile bölünce \(\widehat{\mathbf u}=(3/5,4/5)\) elde edilir. Norm kontrolü \(9/25+16/25=1\) verir. Ters yöndeki birim vektör bunun negatifidir.
\end{enumerate}\end{tcolorbox}
\Needspace{10\baselineskip}\begin{tcolorbox}[title={Örnek 1.2 | Dersteki açı sorusu},colframe=accent]\(\mathbf u=(2,-1,3)\) ve \(\mathbf v=(1,4,2)\) dik midir? \(\cos\theta\) nedir?
\par\vspace{1mm}{\scriptsize\sffamily Ders fotoğrafından | \nolinkurl{IMG-20261009-WA0032.jpg}}\end{tcolorbox}
\begin{tcolorbox}[title={Adım adım çözüm},colback=green!3!white,colframe=green]\begin{enumerate}
\item \(\mathbf u\cdot\mathbf v=2-4+6=4\). Çarpım sıfır olmadığından dik değildirler.
\item \(\|\mathbf u\|=\sqrt{14}\), \(\|\mathbf v\|=\sqrt{21}\). Dolayısıyla \(\cos\theta=4/\sqrt{294}\). Pozitif sonuç açının dar olduğunu doğrular.
\end{enumerate}\end{tcolorbox}
\Needspace{10\baselineskip}\begin{tcolorbox}[title={Örnek 1.3 | İzdüşüm hesabı},colframe=accent]\(\mathbf u=(2,1,2)\) vektörünü \(\mathbf v=(1,0,1)\) doğrultusuna izdüşürünüz.
\par\vspace{1mm}{\scriptsize\sffamily Pekiştirme}\end{tcolorbox}
\begin{tcolorbox}[title={Adım adım çözüm},colback=green!3!white,colframe=green]\begin{enumerate}
\item \(\mathbf u\cdot\mathbf v=4\), \(\|\mathbf v\|^2=2\). Bu yüzden izdüşüm \(2(1,0,1)=(2,0,2)\).
\item Dik kalan \((0,1,0)\) olur ve \((0,1,0)\cdot(1,0,1)=0\). Skaler bileşen \(4/\sqrt2=2\sqrt2\).
\end{enumerate}\end{tcolorbox}
\Needspace{10\baselineskip}\begin{tcolorbox}[title={Örnek 1.4 | Tahtadaki vektörel çarpım},colframe=accent]\(\mathbf u=(1,2,3)\), \(\mathbf v=(2,0,1)\) için \(\mathbf u\times\mathbf v\) bulunuz.
\par\vspace{1mm}{\scriptsize\sffamily Ders fotoğrafından | \nolinkurl{IMG-20261023-WA0029.jpg}}\end{tcolorbox}
\begin{tcolorbox}[title={Adım adım çözüm},colback=green!3!white,colframe=green]\begin{enumerate}
\item Bileşenler \((2\cdot1-3\cdot0,\,3\cdot2-1\cdot1,\,1\cdot0-2\cdot2)\) şeklindedir. Sonuç \((2,5,-4)\).
\item Diklik kontrolleri: \((2,5,-4)\cdot(1,2,3)=2+10-12=0\) ve \((2,5,-4)\cdot(2,0,1)=4-4=0\). Alan \(\sqrt{45}=3\sqrt5\), üçgen alanı \(3\sqrt5/2\).
\end{enumerate}\end{tcolorbox}
\Needspace{10\baselineskip}\begin{tcolorbox}[title={Örnek 1.5 | Hacim ve bağımsızlık},colframe=accent]\(\mathbf u=(1,0,0),\mathbf v=(0,2,0),\mathbf w=(1,1,3)\) için hacim ve bağımsızlık nedir?
\par\vspace{1mm}{\scriptsize\sffamily Pekiştirme}\end{tcolorbox}
\begin{tcolorbox}[title={Adım adım çözüm},colback=green!3!white,colframe=green]\begin{enumerate}
\item \(\mathbf v\times\mathbf w=(6,0,-2)\); \(\mathbf u\cdot(\mathbf v\times\mathbf w)=6\). Paralelyüzlü hacmi 6, dörtyüzlü hacmi 1 birim küptür.
\item Determinant sıfır olmadığından bu üç vektör lineer bağımsızdır ve \(\mathbb R^3\) için baz oluşturur.
\end{enumerate}\end{tcolorbox}
\Needspace{10\baselineskip}\begin{tcolorbox}[title={Örnek 1.6 | Düzlemde bağımlı vektörler},colframe=accent]\(\mathbf u=(1,2)\), \(\mathbf v=(-3,-6)\) bağımsız mıdır?
\par\vspace{1mm}{\scriptsize\sffamily Ders fotoğrafından | \nolinkurl{IMG-20261009-WA0033.jpg}}\end{tcolorbox}
\begin{tcolorbox}[title={Adım adım çözüm},colback=green!3!white,colframe=green]\begin{enumerate}
\item \(\mathbf v=-3\mathbf u\). Bu nedenle \(3\mathbf u+\mathbf v=\mathbf0\), sıfırdan farklı katsayılarla sağlanır.
\item Doğrultuları aynı olduğu için birlikte düzlemi germezler; yalnızca bir doğruyu gererler.
\end{enumerate}\end{tcolorbox}
\Needspace{10\baselineskip}\begin{tcolorbox}[title={Örnek 1.7 | Düzlemde baz},colframe=accent]\(\mathbf u=(2,-3)\), \(\mathbf v=(-5,4)\) düzlem için baz mıdır?
\par\vspace{1mm}{\scriptsize\sffamily Ders fotoğrafından | \nolinkurl{IMG-20261009-WA0034.jpg}}\end{tcolorbox}
\begin{tcolorbox}[title={Adım adım çözüm},colback=green!3!white,colframe=green]\begin{enumerate}
\item Determinant \(2\cdot4-(-3)(-5)=8-15=-7\ne0\). Dolayısıyla bağımsızdırlar ve bazdırlar.
\item Her \((x,y)\), \(\alpha\mathbf u+\beta\mathbf v\) olarak tek biçimde yazılır. Denklem sistemi çözülürse \(\alpha=-(4x+5y)/7\), \(\beta=-(3x+2y)/7\). Yerine koyma her iki koordinatı doğrular.
\end{enumerate}\end{tcolorbox}
\Needspace{10\baselineskip}\begin{tcolorbox}[title={Örnek 1.8 | Pekiştirme},colframe=accent]\(\mathbf a=(1,1,0),\mathbf b=(1,-1,2)\) için diklik, birim vektör ve paralelkenar alanını bulunuz.
\par\vspace{1mm}{\scriptsize\sffamily Pekiştirme}\end{tcolorbox}
\begin{tcolorbox}[title={Adım adım çözüm},colback=green!3!white,colframe=green]\begin{enumerate}
\item \(\mathbf a\cdot\mathbf b=0\), dolayısıyla diktirler. \(\widehat{\mathbf a}=(1,1,0)/\sqrt2\).
\item \(\mathbf a\times\mathbf b=(2,-2,-2)\); alan \(2\sqrt3\). Ayrıca dik vektörlerde alan \(\|\mathbf a\|\|\mathbf b\|=\sqrt2\sqrt6=2\sqrt3\), aynı sonucu verir.
\end{enumerate}\end{tcolorbox}
\clearpage\section{Doğrular, düzlemler ve kuadratik yüzeyler}
\Needspace{10\baselineskip}\begin{tcolorbox}[title={Örnek 2.1 | İki noktadan geçen doğru},colframe=accent]\(A=(0,1,2)\), \(B=(3,1,-1)\) doğrusunu bulunuz.
\par\vspace{1mm}{\scriptsize\sffamily Ders fotoğrafından | \nolinkurl{IMG-20261009-WA0032.jpg}}\end{tcolorbox}
\begin{tcolorbox}[title={Adım adım çözüm},colback=green!3!white,colframe=green]\begin{enumerate}
\item \(\overrightarrow{AB}=(3,0,-3)\), bunun yerine \((1,0,-1)\) kullanılabilir.
\item \(x=t,\ y=1,\ z=2-t\). Kartezyen biçim \(x=2-z,\ y=1\). Parametre 0 iken A, 3 iken B elde edilir.
\end{enumerate}\end{tcolorbox}
\Needspace{10\baselineskip}\begin{tcolorbox}[title={Örnek 2.2 | Nokta ve paralel vektör},colframe=accent]\(A=(1,2,3)\) noktasından geçen ve \((2,3,0)\) yönüne paralel doğruyu yazınız.
\par\vspace{1mm}{\scriptsize\sffamily Ders fotoğrafından | \nolinkurl{IMG-20261009-WA0030.jpg}}\end{tcolorbox}
\begin{tcolorbox}[title={Adım adım çözüm},colback=green!3!white,colframe=green]\begin{enumerate}
\item \(x=1+2t,\ y=2+3t,\ z=3\).
\item Simetrik gösterim \((x-1)/2=(y-2)/3\), ayrıca \(z=3\). İki denklem uzayda tek doğruyu belirler.
\end{enumerate}\end{tcolorbox}
\Needspace{10\baselineskip}\begin{tcolorbox}[title={Örnek 2.3 | Tahtadaki sıfır paydalı gösterim},colframe=accent]Tahtada \(x/1=y/1=z/0\) ve \((x-1)/1=y/(-1)=z/0\) ile belirtilen doğruların açısı nedir?
\par\vspace{1mm}{\scriptsize\sffamily Ders fotoğrafından | \nolinkurl{IMG-20261001-WA0014.jpg}}\end{tcolorbox}
\begin{tcolorbox}[title={Adım adım çözüm},colback=green!3!white,colframe=green]\begin{enumerate}
\item Bu gösterimde \(z/0\) gerçek bir bölüm değildir; \(z=0\) kısıtının kısa yazımı olarak yorumlanır. Parametrik biçimler \((t,t,0)\) ve \((1+s,-s,0)\) olur.
\item Yönler \((1,1,0)\), \((1,-1,0)\), çarpımları 0. \(t=1+s=-s\) den \(t=1/2,s=-1/2\); kesişim \((1/2,1/2,0)\). Dolayısıyla açı \(90^\circ\).
\end{enumerate}\end{tcolorbox}
\Needspace{10\baselineskip}\begin{tcolorbox}[title={Örnek 2.4 | Koordinat eksenlerini kesen düzlem},colframe=accent]\(A=(1,0,0),B=(0,1,0),C=(0,0,1)\) noktalarından geçen düzlemi bulunuz.
\par\vspace{1mm}{\scriptsize\sffamily Ders fotoğrafından | \nolinkurl{IMG-20261001-WA0014.jpg}}\end{tcolorbox}
\begin{tcolorbox}[title={Adım adım çözüm},colback=green!3!white,colframe=green]\begin{enumerate}
\item \(\overrightarrow{AB}=(-1,1,0)\), \(\overrightarrow{AC}=(-1,0,1)\). Çarpımları \((1,1,1)\).
\item \((x-1)+y+z=0\), yani \(x+y+z=1\). Üç nokta da denklemi sağlar.
\end{enumerate}\end{tcolorbox}
\Needspace{10\baselineskip}\begin{tcolorbox}[title={Örnek 2.5 | Nokta ve normal ile düzlem},colframe=accent]\(A=(2,3,4)\) noktasından geçen, normali \((1,3,4)\) olan düzlemi bulunuz.
\par\vspace{1mm}{\scriptsize\sffamily Ders fotoğrafından | \nolinkurl{IMG-20261023-WA0030.jpg}}\end{tcolorbox}
\begin{tcolorbox}[title={Adım adım çözüm},colback=green!3!white,colframe=green]\begin{enumerate}
\item \((x-2)+3(y-3)+4(z-4)=0\).
\item Düzenlenince \(x+3y+4z-27=0\). Normalin herhangi bir sıfır olmayan katı aynı düzlemi verir.
\end{enumerate}\end{tcolorbox}
\Needspace{10\baselineskip}\begin{tcolorbox}[title={Örnek 2.6 | Dersteki doğru-düzlem kesişimi},colframe=accent]\((x-2)/1=y/2=(z+1)/(-1)\) doğrusu ile \(x-y+2z=3\) düzleminin kesişimi nedir?
\par\vspace{1mm}{\scriptsize\sffamily Ders fotoğrafından | \nolinkurl{IMG-20261001-WA0014.jpg}}\end{tcolorbox}
\begin{tcolorbox}[title={Adım adım çözüm},colback=green!3!white,colframe=green]\begin{enumerate}
\item \(x=2+t,\ y=2t,\ z=-1-t\) yazılır.
\item Yerine koyunca \((2+t)-2t+2(-1-t)=3\), yani \(-3t=3\). \(t=-1\) ile nokta \((1,-2,0)\). Kontrol: \(1-(-2)+0=3\).
\end{enumerate}\end{tcolorbox}
\Needspace{10\baselineskip}\begin{tcolorbox}[title={Örnek 2.7 | İki düzlemin arakesiti},colframe=accent]\(x+y+z=2\) ve \(x-y+z=0\) düzlemlerinin kesişimini yazınız.
\par\vspace{1mm}{\scriptsize\sffamily Pekiştirme}\end{tcolorbox}
\begin{tcolorbox}[title={Adım adım çözüm},colback=green!3!white,colframe=green]\begin{enumerate}
\item Denklemleri çıkarınca \(2y=2\), yani \(y=1\). Sonra \(x+z=1\).
\item \(x=t\) seçilirse \(\mathbf r(t)=(t,1,1-t)\). Yön \((1,0,-1)\), normallerin vektörel çarpımına paraleldir.
\end{enumerate}\end{tcolorbox}
\Needspace{10\baselineskip}\begin{tcolorbox}[title={Örnek 2.8 | Tahtadaki nokta-doğru uzaklığı},colframe=accent]\(P=(2,3,-1)\) noktasının \((x-1)/1=(y-2)/1=(z-13)/4\) doğrusuna uzaklığını bulunuz.
\par\vspace{1mm}{\scriptsize\sffamily Ders fotoğrafından | \nolinkurl{IMG-20261023-WA0057.jpg}}\end{tcolorbox}
\begin{tcolorbox}[title={Adım adım çözüm},colback=green!3!white,colframe=green]\begin{enumerate}
\item Doğru üzerinde \(A=(1,2,13)\), yön \(\mathbf u=(1,1,4)\). \(\overrightarrow{AP}=(1,1,-14)\).
\item \(\overrightarrow{AP}\times\mathbf u=(18,-18,0)\); normu \(18\sqrt2\). Yönün normu \(\sqrt{18}=3\sqrt2\). Uzaklık 6.
\end{enumerate}\end{tcolorbox}
\Needspace{10\baselineskip}\begin{tcolorbox}[title={Örnek 2.9 | Çalışma sorusu: nokta-düzlem uzaklığı},colframe=accent]\(P=(1,0,2)\) noktasının \(x+2y-2z+4=0\) düzlemine uzaklığı nedir?
\par\vspace{1mm}{\scriptsize\sffamily Ders fotoğrafından | \nolinkurl{IMG-20261001-WA0014.jpg}}\end{tcolorbox}
\begin{tcolorbox}[title={Adım adım çözüm},colback=green!3!white,colframe=green]\begin{enumerate}
\item Pay \(|1+0-4+4|=1\), payda \(\sqrt{1+4+4}=3\). Uzaklık \(1/3\).
\item Dik izdüşüm noktası \(H=P-\frac{1}{9}(1,2,-2)=(8/9,-2/9,20/9)\). Denkleme koyunca sıfır çıkar ve \(\|P-H\|=1/3\).
\end{enumerate}\end{tcolorbox}
\Needspace{10\baselineskip}\begin{tcolorbox}[title={Örnek 2.10 | Dersteki küre sorusu},colframe=accent]\(x^2+y^2+z^2+4x-6y+2z-11=0\) küresinin merkezi ve yarıçapı nedir?
\par\vspace{1mm}{\scriptsize\sffamily Ders fotoğrafından | \nolinkurl{IMG-20261009-WA0032.jpg}}\end{tcolorbox}
\begin{tcolorbox}[title={Adım adım çözüm},colback=green!3!white,colframe=green]\begin{enumerate}
\item \((x+2)^2-4+(y-3)^2-9+(z+1)^2-1-11=0\).
\item \((x+2)^2+(y-3)^2+(z+1)^2=25\). Merkez \((-2,3,-1)\), yarıçap 5.
\end{enumerate}\end{tcolorbox}
\Needspace{10\baselineskip}\begin{tcolorbox}[title={Örnek 2.11 | Doğruya dik düzlem},colframe=accent]P=(2,3,5) noktasından geçen, \((x-1)/4=(y-2)/(-1)=(z+1)/6\) doğrusuna dik düzlemi bulunuz.
\par\vspace{1mm}{\scriptsize\sffamily Ders fotoğrafından | \nolinkurl{IMG-20261002-WA0024.jpg}}\end{tcolorbox}
\begin{tcolorbox}[title={Adım adım çözüm},colback=green!3!white,colframe=green]\begin{enumerate}
\item Doğrunun yönü \((4,-1,6)\), istenen düzlemin normalidir. \(4(x-2)-(y-3)+6(z-5)=0\).
\item Denklem \(4x-y+6z-35=0\). P yi yerine koyunca sıfır, doğru yönüyle normalin paralelliği ise dikliği doğrular.
\end{enumerate}\end{tcolorbox}
\Needspace{10\baselineskip}\begin{tcolorbox}[title={Örnek 2.12 | Aynı yöntemin ikinci örneği},colframe=accent]A=(2,3,-5) noktasından geçen, normali (7,3,4) olan düzlemi yazınız.
\par\vspace{1mm}{\scriptsize\sffamily Ders fotoğrafından | \nolinkurl{IMG-20261002-WA0020.jpg}}\end{tcolorbox}
\begin{tcolorbox}[title={Adım adım çözüm},colback=green!3!white,colframe=green]\begin{enumerate}
\item \(7(x-2)+3(y-3)+4(z+5)=0\).
\item Sonuç \(7x+3y+4z-3=0\). Üçüncü koordinat -5 olduğu için parantezde z+5 gelir.
\end{enumerate}\end{tcolorbox}
\Needspace{10\baselineskip}\begin{tcolorbox}[title={Örnek 2.13 | Tahtadaki ikinci kesişim örneği},colframe=accent]\((x-1)/2=(y-1)/(-1)=z/3\) doğrusu ile \(2x+3y+z+11=0\) düzleminin kesişimi nedir?
\par\vspace{1mm}{\scriptsize\sffamily Ders fotoğrafından | \nolinkurl{IMG-20261002-WA0026.jpg}}\end{tcolorbox}
\begin{tcolorbox}[title={Adım adım çözüm},colback=green!3!white,colframe=green]\begin{enumerate}
\item \(x=1+2t,y=1-t,z=3t\). Yerine koyunca \(2+4t+3-3t+3t+11=0\), yani \(4t+16=0\).
\item t=-4, kesişim \((-7,5,-12)\). Kontrol \(-14+15-12+11=0\).
\end{enumerate}\end{tcolorbox}
\Needspace{10\baselineskip}\begin{tcolorbox}[title={Örnek 2.14 | İki düzlemin arakesiti: tahta örneği},colframe=accent]\(x+y+z=3\) ve \(x-y+z=1\) düzlemlerinin arakesitini bulunuz.
\par\vspace{1mm}{\scriptsize\sffamily Ders fotoğrafından | \nolinkurl{IMG-20261002-WA0026.jpg}}\end{tcolorbox}
\begin{tcolorbox}[title={Adım adım çözüm},colback=green!3!white,colframe=green]\begin{enumerate}
\item Çıkarınca y=1. z=0 seçilirse x=2; nokta (2,1,0). Normal çarpımı (2,0,-2), eşdeğer yön (1,0,-1).
\item Doğru \((x,y,z)=(2,1,0)+t(1,0,-1)\). Her iki düzlem denklemine de tüm t ler için uyar.
\end{enumerate}\end{tcolorbox}
\Needspace{10\baselineskip}\begin{tcolorbox}[title={Örnek 2.15 | Defterdeki ikinci küre},colframe=accent]\(x^2+y^2+z^2-2x-4y-6z-2=0\) küresinin merkezi ve yarıçapını bulunuz.
\par\vspace{1mm}{\scriptsize\sffamily Ders fotoğrafından | \nolinkurl{IMG-20261009-WA0028.jpg}}\end{tcolorbox}
\begin{tcolorbox}[title={Adım adım çözüm},colback=green!3!white,colframe=green]\begin{enumerate}
\item \((x-1)^2+(y-2)^2+(z-3)^2=16\).
\item Merkez (1,2,3), yarıçap 4. Tam kareler açıldığında sabit 14, sağdan taşınan 16 ile -2 elde edilmesi kontrol sağlar.
\end{enumerate}\end{tcolorbox}
\Needspace{10\baselineskip}\begin{tcolorbox}[title={Örnek 2.16 | İki yüzeyin kesitlerini karşılaştır},colframe=accent]\(x^2+y^2-z^2=1\) ve \(z^2-x^2-y^2=1\) neden farklıdır?
\par\vspace{1mm}{\scriptsize\sffamily Pekiştirme}\end{tcolorbox}
\begin{tcolorbox}[title={Adım adım çözüm},colback=green!3!white,colframe=green]\begin{enumerate}
\item İlkinde \(z=k\) için \(x^2+y^2=1+k^2\). Her k için çember vardır; en dar kesit \(z=0\) düzlemindedir. Tek kanatlı hiperboloid.
\item İkincide \(x^2+y^2=k^2-1\). \(|k|<1\) için kesit yok; \(k=\pm1\) için bir nokta, dışarıda çember vardır. Çift kanatlı hiperboloid.
\end{enumerate}\end{tcolorbox}
\Needspace{10\baselineskip}\begin{tcolorbox}[title={Örnek 2.17 | Pekiştirme: yüzeyi tanı},colframe=accent]\(z=4-x^2-y^2\) yüzeyinin tepesini, açılma yönünü ve \(z=0\) kesitini bulunuz.
\par\vspace{1mm}{\scriptsize\sffamily Pekiştirme}\end{tcolorbox}
\begin{tcolorbox}[title={Adım adım çözüm},colback=green!3!white,colframe=green]\begin{enumerate}
\item Tepe \((0,0,4)\), karelerin katsayıları negatif olduğundan aşağı açılan eliptik paraboloid.
\item \(z=0\) kesiti \(x^2+y^2=4\), yarıçapı 2 olan çember. \(z>4\) için kesit yoktur.
\end{enumerate}\end{tcolorbox}
\clearpage\section{Vektör fonksiyonları, eğriler ve hareket}
\Needspace{10\baselineskip}\begin{tcolorbox}[title={Örnek 3.1 | Defterdeki tanım kümesi örneği},colframe=accent]\(\mathbf F(t)=(t+1,\ln(1-t),\sqrt{4-t^2})\) için tanım kümesini bulunuz.
\par\vspace{1mm}{\scriptsize\sffamily Ders fotoğrafından | \nolinkurl{IMG-20261009-WA0016.jpg}}\end{tcolorbox}
\begin{tcolorbox}[title={Adım adım çözüm},colback=green!3!white,colframe=green]\begin{enumerate}
\item İlk bileşen tüm gerçek sayılarda tanımlıdır. Logaritma için \(1-t>0\), yani \(t<1\). Karekök için \(-2\le t\le2\).
\item Koşulları birlikte sağlamak gerekir: \(I=[-2,1)\). \(-2\) kabul edilir, çünkü kök sıfır olabilir; 1 kabul edilmez, çünkü logaritma argümanı sıfır olamaz.
\end{enumerate}\end{tcolorbox}
\Needspace{10\baselineskip}\begin{tcolorbox}[title={Örnek 3.2 | Vektör fonksiyonlarıyla işlemler},colframe=accent]\(\mathbf F=(t,t^2,t^3)\), \(\mathbf G=(1-t,1+t^2,t)\). Toplamı, skaler ve vektörel çarpımı bulunuz.
\par\vspace{1mm}{\scriptsize\sffamily Ders fotoğrafından | \nolinkurl{IMG-20261009-WA0017.jpg}}\end{tcolorbox}
\begin{tcolorbox}[title={Adım adım çözüm},colback=green!3!white,colframe=green]\begin{enumerate}
\item \(\mathbf F+\mathbf G=(1,1+2t^2,t^3+t)\).
\item \(\mathbf F\cdot\mathbf G=t(1-t)+t^2(1+t^2)+t^4=t+2t^4\).
\item \(\mathbf F\times\mathbf G=(-t^5,t^3-t^4-t^2,t+2t^3-t^2)\). Her bileşen determinant açılımından gelir.
\end{enumerate}\end{tcolorbox}
\Needspace{10\baselineskip}\begin{tcolorbox}[title={Örnek 3.3 | Defterdeki limit işlemleri},colframe=accent]\(\mathbf F(t)=(\sin t,\cos t,1+t)\), \(\mathbf G(t)=(\cos2t,\sin3t,e^t+1)\) için t sıfıra giderken skaler ve vektörel çarpım limitlerini bulunuz.
\par\vspace{1mm}{\scriptsize\sffamily Ders fotoğrafından | \nolinkurl{IMG-20261009-WA0019.jpg}}\end{tcolorbox}
\begin{tcolorbox}[title={Adım adım çözüm},colback=green!3!white,colframe=green]\begin{enumerate}
\item \(\mathbf F\to(0,1,1)\), \(\mathbf G\to(1,0,2)\).
\item Skaler çarpım limiti \((0,1,1)\cdot(1,0,2)=2\). Vektörel çarpım limiti \((0,1,1)\times(1,0,2)=(2,1,-1)\).
\end{enumerate}\end{tcolorbox}
\Needspace{10\baselineskip}\begin{tcolorbox}[title={Örnek 3.4 | Tam kısım içeren süreklilik},colframe=accent]\(\mathbf F(t)=(|t|+1,t\lfloor t\rfloor,2+t^2)\) fonksiyonu 0 ve 1 noktalarında sürekli midir?
\par\vspace{1mm}{\scriptsize\sffamily Ders fotoğrafından | \nolinkurl{IMG-20261009-WA0019.jpg}}\end{tcolorbox}
\begin{tcolorbox}[title={Adım adım çözüm},colback=green!3!white,colframe=green]\begin{enumerate}
\item İlk ve üçüncü bileşen her yerde süreklidir. İkinci bileşen 0 yakınında solda \(-t\), sağda 0 olur; iki yandan limit 0, değer de 0. Dolayısıyla 0 noktasında süreklidir.
\item 1 yakınında solda ikinci bileşen 0, sağda t olur. Limitler 0 ve 1 olduğundan 1 noktasında süreksizdir.
\end{enumerate}\end{tcolorbox}
\Needspace{10\baselineskip}\begin{tcolorbox}[title={Örnek 3.5 | Çember ve helis},colframe=accent]\(\mathbf r(t)=(a\cos t,a\sin t,0)\) ile \(\mathbf h(t)=(\cos t,\sin t,3t)\) eğrilerini yorumlayınız.
\par\vspace{1mm}{\scriptsize\sffamily Ders fotoğrafından | \nolinkurl{IMG-20261009-WA0022.jpg}}\end{tcolorbox}
\begin{tcolorbox}[title={Adım adım çözüm},colback=green!3!white,colframe=green]\begin{enumerate}
\item İlkinde \(x^2+y^2=a^2\), \(z=0\). \(a>0\) ve \(0\le t\le2\pi\) için pozitif z yönünden bakıldığında çember saat yönünün tersine bir tur taranır.
\item İkincide \(x^2+y^2=1\); yükseklik t ile artar. Bir turda t, \(2\pi\) artar ve yükseklik \(6\pi\) değişir. Bu eğri bir helistir.
\end{enumerate}\end{tcolorbox}
\Needspace{10\baselineskip}\begin{tcolorbox}[title={Örnek 3.6 | Parçalı düzgün eğri},colframe=accent]\(\mathbf r(t)=(a\cos t,a|\sin t|,0)\), \(-2\pi\le t\le2\pi\), \(a>0\) için düzgünlük nasıldır?
\par\vspace{1mm}{\scriptsize\sffamily Ders fotoğrafından | \nolinkurl{IMG-20261015-WA0038.jpg}}\end{tcolorbox}
\begin{tcolorbox}[title={Adım adım çözüm},colback=green!3!white,colframe=green]\begin{enumerate}
\item \(\sin t\) işaret değiştirdiğinde mutlak değerin türevi değişir. \(t=k\pi\) noktalarında tek taraflı teğet vektörler uyuşmaz; verilen parametrizasyon düzgün değildir.
\item Aradaki açık aralıklarda türev vardır ve normu a olur. Dolayısıyla sonlu sayıda parçaya ayrılarak parçalı düzgün kabul edilir. İz üst yarım çemberdir; birden fazla kez taranır.
\end{enumerate}\end{tcolorbox}
\Needspace{10\baselineskip}\begin{tcolorbox}[title={Örnek 3.7 | Dersteki temel türevler},colframe=accent]\(\mathbf F(t)=(t,t^2,t^3)\) ve \(\mathbf G(t)=(t\ln t-t,e^t,e^{-t})\) türevlerini bulunuz.
\par\vspace{1mm}{\scriptsize\sffamily Ders fotoğrafından | \nolinkurl{IMG-20261023-WA0058.jpg}}\end{tcolorbox}
\begin{tcolorbox}[title={Adım adım çözüm},colback=green!3!white,colframe=green]\begin{enumerate}
\item Bileşenler ayrı türevlenir: \(\mathbf F\prime=(1,2t,3t^2)\).
\item İkinci fonksiyonun ilk bileşeninde çarpım kuralı \((t\ln t-t)\prime=\ln t+1-1=\ln t\) verir. \(\mathbf G\prime=(\ln t,e^t,-e^{-t})\), \(t>0\).
\end{enumerate}\end{tcolorbox}
\Needspace{10\baselineskip}\begin{tcolorbox}[title={Örnek 3.8 | Zincir kuralı içeren vektör türevi},colframe=accent]\(\mathbf F(t)=(\arcsin(4t),-3\arctan(2t-1),7\ln(3t^2))\) türevini bulunuz.
\par\vspace{1mm}{\scriptsize\sffamily Ders fotoğrafından | \nolinkurl{IMG-20261023-WA0058.jpg}}\end{tcolorbox}
\begin{tcolorbox}[title={Adım adım çözüm},colback=green!3!white,colframe=green]\begin{enumerate}
\item İlk bileşenin türevi \(4/\sqrt{1-16t^2}\); ikincinin türevi \(-6/[1+(2t-1)^2]\).
\item Üçüncü bileşen \(7(6t)/(3t^2)=14/t\). Türev formüllerinin ortak geçerlilik aralığı \(-1/4<t<1/4\), \(t\ne0\). Fonksiyon uçlarda tanımlı olsa da bu türevler uçlarda sonlu değildir.
\end{enumerate}\end{tcolorbox}
\Needspace{10\baselineskip}\begin{tcolorbox}[title={Örnek 3.9 | Tahtadaki özdeşliğin kanıtı},colframe=accent]\([\mathbf F(t)\times\mathbf F\prime(t)]\prime=\mathbf F(t)\times\mathbf F\prime\prime(t)\) olduğunu gösteriniz.
\par\vspace{1mm}{\scriptsize\sffamily Ders fotoğrafından | \nolinkurl{IMG-20261023-WA0058.jpg}}\end{tcolorbox}
\begin{tcolorbox}[title={Adım adım çözüm},colback=green!3!white,colframe=green]\begin{enumerate}
\item Çarpım kuralı sol tarafı \(\mathbf F\prime\times\mathbf F\prime+\mathbf F\times\mathbf F\prime\prime\) yapar.
\item Bir vektörün kendisiyle vektörel çarpımı sıfırdır. İlk terim silinince istenen özdeşlik kalır.
\end{enumerate}\end{tcolorbox}
\Needspace{10\baselineskip}\begin{tcolorbox}[title={Örnek 3.10 | Başlangıç koşullu integral},colframe=accent]\(\mathbf r\prime(t)=(2t,e^{3t},4\sin2t)\), \(\mathbf r(0)=(2,1,-2)\). Konumu bulunuz.
\par\vspace{1mm}{\scriptsize\sffamily Ders fotoğrafından | \nolinkurl{IMG-20261015-WA0037.jpg}}\end{tcolorbox}
\begin{tcolorbox}[title={Adım adım çözüm},colback=green!3!white,colframe=green]\begin{enumerate}
\item İntegral \(\mathbf r(t)=(t^2+C_1,e^{3t}/3+C_2,-2\cos2t+C_3)\).
\item Başlangıçtan \(C_1=2,C_2=2/3,C_3=0\). Dolayısıyla \(\mathbf r(t)=(t^2+2,(e^{3t}+2)/3,-2\cos2t)\).
\end{enumerate}\end{tcolorbox}
\Needspace{10\baselineskip}\begin{tcolorbox}[title={Örnek 3.11 | Dersteki belirli integral üçlüsü},colframe=accent]\(\int_0^1(4t^3,3,2(t+1))dt\), \(\int_1^2(6-6t,3\sqrt t,4/t^2)dt\) ve \(\int_{-\pi/4}^{\pi/4}(\sin t,1+\cos t,\sec^2t)dt\) hesaplayınız.
\par\vspace{1mm}{\scriptsize\sffamily Ders fotoğrafından | \nolinkurl{IMG-20261023-WA0058.jpg}}\end{tcolorbox}
\begin{tcolorbox}[title={Adım adım çözüm},colback=green!3!white,colframe=green]\begin{enumerate}
\item İlkinde ilkel fonksiyon \((t^4,3t,t^2+2t)\). Sonuç \((1,3,3)\).
\item İkincide ilkel fonksiyon \((6t-3t^2,2t^{3/2},-4/t)\). Sonuç \((-3,4\sqrt2-2,2)\).
\item Üçüncüde sinüs tek fonksiyon olduğundan ilk bileşen 0. İkinci \(\pi/2+\sqrt2\), üçüncü \(\tan(\pi/4)-\tan(-\pi/4)=2\).
\end{enumerate}\end{tcolorbox}
\Needspace{10\baselineskip}\begin{tcolorbox}[title={Örnek 3.12 | Dersteki 5-12-13 helisi},colframe=accent]\(\mathbf r(t)=(5\cos t,5\sin t,12t)\), \(0\le t\le2\pi\), uzunluğunu bulunuz.
\par\vspace{1mm}{\scriptsize\sffamily Ders fotoğrafından | \nolinkurl{IMG-20261015-WA0039.jpg}}\end{tcolorbox}
\begin{tcolorbox}[title={Adım adım çözüm},colback=green!3!white,colframe=green]\begin{enumerate}
\item \(\mathbf r\prime=(-5\sin t,5\cos t,12)\); norm \(\sqrt{25\sin^2t+25\cos^2t+144}=13\).
\item \(L=\int_0^{2\pi}13dt=26\pi\). Bu aralıkta \(s=13t\) olduğundan yay parametreli biçim \((5\cos(s/13),5\sin(s/13),12s/13)\), \(0\le s\le26\pi\).
\end{enumerate}\end{tcolorbox}
\Needspace{10\baselineskip}\begin{tcolorbox}[title={Örnek 3.13 | Çemberin Frenet çatısı},colframe=accent]\(\mathbf r(t)=(\cos t,\sin t,0)\) için \(\mathbf T,\mathbf N,\mathbf B\) nedir?
\par\vspace{1mm}{\scriptsize\sffamily Ders fotoğrafından | \nolinkurl{IMG-20261015-WA0034.jpg}}\end{tcolorbox}
\begin{tcolorbox}[title={Adım adım çözüm},colback=green!3!white,colframe=green]\begin{enumerate}
\item \(\mathbf r\prime=(-\sin t,\cos t,0)\) normu 1 olduğundan \(\mathbf T=(-\sin t,\cos t,0)\).
\item \(\mathbf T\prime=(-\cos t,-\sin t,0)\) normu da 1: \(\mathbf N=(-\cos t,-\sin t,0)\).
\item Çarpım \(\mathbf B=(0,0,1)\). Normal merkeze yönelir; binormal sabittir.
\end{enumerate}\end{tcolorbox}
\Needspace{10\baselineskip}\begin{tcolorbox}[title={Örnek 3.14 | Dersteki uzay eğrisinin Frenet üçlüsü},colframe=accent]\(\mathbf r(t)=(e^t,e^{-t},\sqrt2\,t)\) için \(t=0\) anındaki çatıyı bulunuz.
\par\vspace{1mm}{\scriptsize\sffamily Ders fotoğrafından | \nolinkurl{IMG-20261015-WA0040.jpg}}\end{tcolorbox}
\begin{tcolorbox}[title={Adım adım çözüm},colback=green!3!white,colframe=green]\begin{enumerate}
\item \(\mathbf v(0)=(1,-1,\sqrt2)\), normu 2; \(\mathbf a(0)=(1,1,0)\). Dolayısıyla \(\mathbf T(0)=(1,-1,\sqrt2)/2\).
\item \(\mathbf v(0)\cdot\mathbf a(0)=0\), süratin türevi o anda 0. \(\mathbf T\prime(0)=\mathbf a(0)/2\); normu \(1/\sqrt2\). Böylece \(\mathbf N(0)=(1,1,0)/\sqrt2\).
\item \(\mathbf B(0)=(-1/2,1/2,1/\sqrt2)\). Üç vektörün normları 1 ve ikili skaler çarpımları sıfırdır.
\end{enumerate}\end{tcolorbox}
\Needspace{10\baselineskip}\begin{tcolorbox}[title={Örnek 3.15 | Genel helisin eğriliği},colframe=accent]\(\mathbf r(t)=(a\cos t,a\sin t,bt)\), \(a>0\), için eğrilik nedir?
\par\vspace{1mm}{\scriptsize\sffamily Pekiştirme}\end{tcolorbox}
\begin{tcolorbox}[title={Adım adım çözüm},colback=green!3!white,colframe=green]\begin{enumerate}
\item \(\|\mathbf r\prime\|=\sqrt{a^2+b^2}\). \(\mathbf r\prime\times\mathbf r\prime\prime=(ab\sin t,-ab\cos t,a^2)\), normu \(a\sqrt{a^2+b^2}\).
\item \(\kappa=a/(a^2+b^2)\), \(\rho=(a^2+b^2)/a\). \(b=0\) olursa çember sonucu \(1/a\) elde edilir.
\end{enumerate}\end{tcolorbox}
\Needspace{10\baselineskip}\begin{tcolorbox}[title={Örnek 3.16 | Tahtadaki hareket örneği},colframe=accent]\(\mathbf r(t)=(5\cos t,5\sin t,6t^2)\). Hız, ivme ve bunların dik olduğu anı bulunuz.
\par\vspace{1mm}{\scriptsize\sffamily Ders fotoğrafından | \nolinkurl{IMG-20261015-WA0036.jpg}}\end{tcolorbox}
\begin{tcolorbox}[title={Adım adım çözüm},colback=green!3!white,colframe=green]\begin{enumerate}
\item \(\mathbf v=(-5\sin t,5\cos t,12t)\), \(\mathbf a=(-5\cos t,-5\sin t,12)\).
\item \(\mathbf v\cdot\mathbf a=25\sin t\cos t-25\sin t\cos t+144t=144t\). Yalnızca \(t=0\) anında diktirler.
\end{enumerate}\end{tcolorbox}
\Needspace{10\baselineskip}\begin{tcolorbox}[title={Örnek 3.17 | Üstel ve trigonometrik hareket},colframe=accent]\(\mathbf r(t)=(e^{-t},2\cos3t,2\sin3t)\). \(t=0\) anında hız, ivme büyüklükleri ve birim teğeti bulunuz.
\par\vspace{1mm}{\scriptsize\sffamily Ders fotoğrafından | \nolinkurl{IMG-20261015-WA0036.jpg}}\end{tcolorbox}
\begin{tcolorbox}[title={Adım adım çözüm},colback=green!3!white,colframe=green]\begin{enumerate}
\item \(\mathbf v=(-e^{-t},-6\sin3t,6\cos3t)\), \(\mathbf a=(e^{-t},-18\cos3t,-18\sin3t)\).
\item \(\mathbf v(0)=(-1,0,6)\), \(\|\mathbf v(0)\|=\sqrt{37}\); \(\mathbf a(0)=(1,-18,0)\), \(\|\mathbf a(0)\|=\sqrt{325}=5\sqrt{13}\).
\item \(\mathbf T(0)=(-1,0,6)/\sqrt{37}\). Vektörü büyüklüğüyle karıştırmamak için sonucun türünü kontrol et.
\end{enumerate}\end{tcolorbox}
\Needspace{10\baselineskip}\begin{tcolorbox}[title={Örnek 3.18 | Parametrik eğrinin türü},colframe=accent]\(\mathbf r(t)=(2+3t,4+5t,6+4t)\) eğrisinin türünü bulunuz.
\par\vspace{1mm}{\scriptsize\sffamily Ders fotoğrafından | \nolinkurl{IMG-20261009-WA0021.jpg}}\end{tcolorbox}
\begin{tcolorbox}[title={Adım adım çözüm},colback=green!3!white,colframe=green]\begin{enumerate}
\item \(\mathbf r(t)=(2,4,6)+t(3,5,4)\); sabit doğrultuda bir doğrudur.
\item Simetrik biçim \((x-2)/3=(y-4)/5=(z-6)/4\).
\end{enumerate}\end{tcolorbox}
\Needspace{10\baselineskip}\begin{tcolorbox}[title={Örnek 3.19 | Kapalı eğri ve uç teğetler},colframe=accent]\(\mathbf r(t)=(t^2-4,t^3-4t,0)\), \(-2\le t\le2\). Basit ve kapalı mıdır?
\par\vspace{1mm}{\scriptsize\sffamily Ders fotoğrafından | \nolinkurl{IMG-20261009-WA0025.jpg}}\end{tcolorbox}
\begin{tcolorbox}[title={Adım adım çözüm},colback=green!3!white,colframe=green]\begin{enumerate}
\item Her iki uç (0,0,0) olur. Farklı t ve s için eşit ilk koordinatlar s=-t gerektirir; ikinci koordinat eşitliği \(2t(t^2-4)=0\) verir.
\item Farklı parametrelerle tek tekrar uçlarda t=2 ve s=-2 dir; t=0 kendi kendine aynı parametredir. Dolayısıyla basit kapalı eğridir.
\item \(\mathbf r\prime=(-4,8,0)\) ve (4,8,0) uçlarda paralel değildir; kapama noktasında köşe vardır. Verilen parametrizasyon parçalı düzgün kapalı eğri olarak ele alınır.
\end{enumerate}\end{tcolorbox}
\Needspace{10\baselineskip}\begin{tcolorbox}[title={Örnek 3.20 | Türev ve ikinci türevin noktadaki değeri},colframe=accent]\(\mathbf r(t)=(1-\cos t,\sin t,t)\) için \(t=\pi/2\) anındaki ilk iki türevi bulunuz.
\par\vspace{1mm}{\scriptsize\sffamily Ders fotoğrafından | \nolinkurl{IMG-20261009-WA0024.jpg}}\end{tcolorbox}
\begin{tcolorbox}[title={Adım adım çözüm},colback=green!3!white,colframe=green]\begin{enumerate}
\item \(\mathbf r\prime=(\sin t,\cos t,1)\), \(\mathbf r\prime\prime=(\cos t,-\sin t,0)\).
\item Noktadaki değerler (1,0,1) ve (0,-1,0). Teğet noktası \((1,1,\pi/2)\), teğet doğrusu bu noktadan (1,0,1) yönünde geçer.
\end{enumerate}\end{tcolorbox}
\Needspace{10\baselineskip}\begin{tcolorbox}[title={Örnek 3.21 | Vektörel çarpımın türevini iki yoldan almak},colframe=accent]\(\mathbf U=(t,-t,t^2/2)\), \(\mathbf V=(t^4/4,t^3/3,t^2/2)\). Çarpımın türevini bulunuz.
\par\vspace{1mm}{\scriptsize\sffamily Ders fotoğrafından | \nolinkurl{IMG-20261009-WA0024.jpg}}\end{tcolorbox}
\begin{tcolorbox}[title={Adım adım çözüm},colback=green!3!white,colframe=green]\begin{enumerate}
\item Önce çarpım: \(\mathbf U\times\mathbf V=(-t^3/2-t^5/6,t^6/8-t^3/2,t^4/3+t^5/4)\).
\item Türev \((-3t^2/2-5t^4/6,3t^5/4-3t^2/2,4t^3/3+5t^4/4)\).
\item Alternatif \(\mathbf U\prime\times\mathbf V+\mathbf U\times\mathbf V\prime\) aynı sonucu verir. Orta bileşenin t altıncı kuvvetinin türevi \(6t^5/8=3t^5/4\); defterdeki bu satır sadeleştirilerek kontrol edildi.
\end{enumerate}\end{tcolorbox}
\Needspace{10\baselineskip}\begin{tcolorbox}[title={Örnek 3.22 | Sabit normun türevi},colframe=accent]\(\mathbf r(t)=(a\cos t,a\sin t,0)\), a>0, için hızın konuma dik olduğunu kanıtlayınız.
\par\vspace{1mm}{\scriptsize\sffamily Ders fotoğrafından | \nolinkurl{IMG-20261015-WA0036.jpg}}\end{tcolorbox}
\begin{tcolorbox}[title={Adım adım çözüm},colback=green!3!white,colframe=green]\begin{enumerate}
\item \(\mathbf r\prime=(-a\sin t,a\cos t,0)\); konumla çarpım \(-a^2\sin t\cos t+a^2\sin t\cos t=0\).
\item Genel gerekçe: \(\mathbf r\cdot\mathbf r=a^2\) sabit olduğundan türevi \(2\mathbf r\cdot\mathbf r\prime=0\).
\end{enumerate}\end{tcolorbox}
\Needspace{10\baselineskip}\begin{tcolorbox}[title={Örnek 3.23 | Dairesel harekette ivme},colframe=accent]\(\mathbf r(t)=(a\cos\omega t,a\sin\omega t,0)\), \(a>0,\omega\ne0\) için \(a_T,a_N\) nedir?
\par\vspace{1mm}{\scriptsize\sffamily Pekiştirme}\end{tcolorbox}
\begin{tcolorbox}[title={Adım adım çözüm},colback=green!3!white,colframe=green]\begin{enumerate}
\item Sürat \(a|\omega|\) sabit olduğundan \(a_T=0\). Eğrilik \(1/a\).
\item \(a_N=(1/a)(a^2\omega^2)=a\omega^2\). İvme tamamen içeriye yönelir; sürat sabit olmasına rağmen sıfır değildir.
\end{enumerate}\end{tcolorbox}
\clearpage\section{Çok değişkenli fonksiyonlar, limit ve süreklilik}
\Needspace{10\baselineskip}\begin{tcolorbox}[title={Örnek 4.1 | Disk ve eğik şerit},colframe=accent]\(f(x,y)=\sqrt{64-x^2-y^2}\) ve \(g(x,y)=\arccos(x-y)\) için tanım ve görüntü kümelerini bulunuz.
\par\vspace{1mm}{\scriptsize\sffamily Ders fotoğrafından | \nolinkurl{IMG-20261015-WA0041.jpg}}\end{tcolorbox}
\begin{tcolorbox}[title={Adım adım çözüm},colback=green!3!white,colframe=green]\begin{enumerate}
\item İlkinde \(x^2+y^2\le64\): merkezi orijin, yarıçapı 8 olan kapalı disk. Çıktılar \([0,8]\); grafik üst yarım küredir.
\item İkincide \(-1\le x-y\le1\), yani \(x-1\le y\le x+1\). Sınır doğruları dahil bir şerit. Görüntü \([0,\pi]\).
\end{enumerate}\end{tcolorbox}
\Needspace{10\baselineskip}\begin{tcolorbox}[title={Örnek 4.2 | Kök ve logaritmayı birlikte ele almak},colframe=accent]\(f(x,y)=\sqrt{xy}+\log(1-x^2-y^2)\) fonksiyonunun tanım bölgesini bulunuz.
\par\vspace{1mm}{\scriptsize\sffamily Ders fotoğrafından | \nolinkurl{IMG-20261015-WA0042.jpg}}\end{tcolorbox}
\begin{tcolorbox}[title={Adım adım çözüm},colback=green!3!white,colframe=green]\begin{enumerate}
\item \(xy\ge0\) ve \(x^2+y^2<1\) birlikte gerekli.
\item Bölge açık birim diskin birinci ve üçüncü bölgelerde kalan parçaları ile disk içindeki eksen parçalarıdır. Eksenler dahildir; çember dahil değildir.
\end{enumerate}\end{tcolorbox}
\Needspace{10\baselineskip}\begin{tcolorbox}[title={Örnek 4.3 | İşaret ayırımı gerektiren kök},colframe=accent]\(f(x,y)=\sqrt{xy-x^4}+xy\) ve \(g(x,y)=\sqrt{x^2-y}+x^2+y^2+1\) tanım bölgelerini bulunuz.
\par\vspace{1mm}{\scriptsize\sffamily Ders fotoğrafından | \nolinkurl{IMG-20261022-WA0053.jpg}}\end{tcolorbox}
\begin{tcolorbox}[title={Adım adım çözüm},colback=green!3!white,colframe=green]\begin{enumerate}
\item İlkinde \(x(y-x^3)\ge0\). \(x>0\) iken \(y\ge x^3\), \(x<0\) iken \(y\le x^3\); \(x=0\) iken her y geçerlidir. Sonuç \(\{x>0,y\ge x^3\}\cup\{x<0,y\le x^3\}\cup\{x=0\}\).
\item İkincide \(y\le x^2\). Parabol ve altı dahildir. İlk soruda x ile bölmeden önce x işareti ayrılmazsa bölgenin yarısı kaybolur.
\end{enumerate}\end{tcolorbox}
\Needspace{10\baselineskip}\begin{tcolorbox}[title={Örnek 4.4 | Her yerde tanımlı logaritma},colframe=accent]\(\ln(x^2+y^2+4)\) için tanım kümesi nedir?
\par\vspace{1mm}{\scriptsize\sffamily Ders fotoğrafından | \nolinkurl{IMG-20261022-WA0053.jpg}}\end{tcolorbox}
\begin{tcolorbox}[title={Adım adım çözüm},colback=green!3!white,colframe=green]\begin{enumerate}
\item \(x^2+y^2+4\ge4>0\) her noktada doğrudur. Dolayısıyla tanım kümesi \(\mathbb R^2\).
\item Logaritma artan olduğundan görüntü kümesi \([\ln4,\infty)\).
\end{enumerate}\end{tcolorbox}
\Needspace{10\baselineskip}\begin{tcolorbox}[title={Örnek 4.5 | Paraboloidin seviye eğrileri},colframe=accent]\(f(x,y)=x^2+y^2\) için \(c=0,1,4\) seviyelerini bulunuz.
\par\vspace{1mm}{\scriptsize\sffamily Ders fotoğrafından | \nolinkurl{IMG-20261015-WA0042.jpg}}\end{tcolorbox}
\begin{tcolorbox}[title={Adım adım çözüm},colback=green!3!white,colframe=green]\begin{enumerate}
\item \(c=0\) için yalnızca orijin, \(c=1\) için yarıçapı 1 çember, \(c=4\) için yarıçapı 2 çember vardır.
\item \(c<0\) için seviye kümesi boştur. Yüzey \(z=x^2+y^2\) yukarı açılan paraboloiddir.
\end{enumerate}\end{tcolorbox}
\Needspace{10\baselineskip}\begin{tcolorbox}[title={Örnek 4.6 | Polinom limiti},colframe=accent]\(P(x,y)=5x^2y^2+4xy^3-8xy^2+3\) için \((x,y)\to(1,2)\) limitini bulunuz.
\par\vspace{1mm}{\scriptsize\sffamily Ders fotoğrafından | \nolinkurl{IMG-20261022-WA0053.jpg}}\end{tcolorbox}
\begin{tcolorbox}[title={Adım adım çözüm},colback=green!3!white,colframe=green]\begin{enumerate}
\item Polinom her yerde süreklidir; limit \(P(1,2)\).
\item \(5\cdot1\cdot4+4\cdot1\cdot8-8\cdot1\cdot4+3=20+32-32+3=23\).
\end{enumerate}\end{tcolorbox}
\Needspace{10\baselineskip}\begin{tcolorbox}[title={Örnek 4.7 | Dersin kısa limitleri},colframe=accent]\(\lim_{(x,y)\to(0,0)}(2x^2+y^2+1)/(x^2+y^2+3)\), \(\lim_{(x,y)\to(0,\ln2)}e^{x-y}\), \(\lim_{(x,y)\to(1,0)}x\sin y/(x^2+1)\) nedir?
\par\vspace{1mm}{\scriptsize\sffamily Ders fotoğrafından | \nolinkurl{IMG-20261022-WA0053.jpg}}\end{tcolorbox}
\begin{tcolorbox}[title={Adım adım çözüm},colback=green!3!white,colframe=green]\begin{enumerate}
\item İlk payda 3 olur; limit \(1/3\).
\item İkinci ifade süreklidir: \(e^{-\ln2}=1/2\).
\item Üçüncü payda 2 ve sinüs sıfır olur; limit 0.
\end{enumerate}\end{tcolorbox}
\Needspace{10\baselineskip}\begin{tcolorbox}[title={Örnek 4.8 | Eşlenikle sadeleştirme},colframe=accent]\(\lim_{(x,y)\to(0,0)}[2-\sqrt{xy+4}]/(xy)\) hesaplayınız.
\par\vspace{1mm}{\scriptsize\sffamily Ders fotoğrafından | \nolinkurl{IMG-20261022-WA0053.jpg}}\end{tcolorbox}
\begin{tcolorbox}[title={Adım adım çözüm},colback=green!3!white,colframe=green]\begin{enumerate}
\item İfade kendi tanım kümesinde \(xy\ne0\) koşuluyla incelenir. Eşlenikle çarpınca pay \(4-(xy+4)=-xy\).
\item Yakında \(-1/[2+\sqrt{xy+4}]\) ile aynı olur. Bu sürekli ifadenin limiti \(-1/4\). Fonksiyonun eksenlerde tanımsız olması limit almayı engellemez.
\end{enumerate}\end{tcolorbox}
\Needspace{10\baselineskip}\begin{tcolorbox}[title={Örnek 4.9 | Sinüs içeren limitler},colframe=accent]\(\lim_{(x,y)\to(0,0)}\sin(xy)/x\) ve \(\lim_{(x,y)\to(0,0)}e^y\sin x/x\) hesaplayınız.
\par\vspace{1mm}{\scriptsize\sffamily Ders fotoğrafından | \nolinkurl{IMG-20261022-WA0056.jpg}}\end{tcolorbox}
\begin{tcolorbox}[title={Adım adım çözüm},colback=green!3!white,colframe=green]\begin{enumerate}
\item İlk ifade, uygun noktalarda \(y[\sin(xy)/(xy)]\) olur. Standart \(\sin u/u\to1\) limiti ve \(y\to0\) ile sonuç 0. Doğrudan \(|\sin(xy)/x|\le|y|\) sınırı da aynı sonucu verir.
\item İkincide \(e^y\to1\) ve \(\sin x/x\to1\); sonuç 1. Her iki limitte yaklaşım fonksiyonun tanımlı olduğu noktalar üzerindendir.
\end{enumerate}\end{tcolorbox}
\Needspace{10\baselineskip}\begin{tcolorbox}[title={Örnek 4.10 | Doğruya bağlı sonuç},colframe=accent]\(f(x,y)=2xy/(x^2+y^2)\) için orijindeki limiti inceleyiniz.
\par\vspace{1mm}{\scriptsize\sffamily Ders fotoğrafından | \nolinkurl{IMG-20261022-WA0058.jpg}}\end{tcolorbox}
\begin{tcolorbox}[title={Adım adım çözüm},colback=green!3!white,colframe=green]\begin{enumerate}
\item \(y=0\) yolunda değer 0; \(y=x\) yolunda, \(x\ne0\) için değer 1.
\item İki yol aynı noktaya gider fakat sonuçlar farklıdır. Dolayısıyla limit yoktur. Genel \(y=mx\) yolunda \(2m/(1+m^2)\) çıkar.
\end{enumerate}\end{tcolorbox}
\Needspace{10\baselineskip}\begin{tcolorbox}[title={Örnek 4.11 | Parabol yoluyla limitin yokluğu},colframe=accent]\(f(x,y)=(x^4-y^2)/(x^4+y^2)\) için orijindeki limiti inceleyiniz.
\par\vspace{1mm}{\scriptsize\sffamily Ders fotoğrafından | \nolinkurl{IMG-20261022-WA0056.jpg}}\end{tcolorbox}
\begin{tcolorbox}[title={Adım adım çözüm},colback=green!3!white,colframe=green]\begin{enumerate}
\item \(y=0\) boyunca değer 1; \(y=x^2\) boyunca değer 0. İki sonuç farklı, limit yok.
\item Genel \(y=ax^2\) yolunda \((1-a^2)/(1+a^2)\) çıkar. Bu örnek, paydadaki dereceleri eşleştirerek uygun yol seçmeyi öğretir.
\end{enumerate}\end{tcolorbox}
\Needspace{10\baselineskip}\begin{tcolorbox}[title={Örnek 4.12 | İki başka karşı örnek},colframe=accent]\((x-y)/(x+y)\), \(x/\sqrt{x^2+y^2}\) ve \((x^2+y^2)/y^2\) orijinde limite sahip midir?
\par\vspace{1mm}{\scriptsize\sffamily Ders fotoğrafından | \nolinkurl{IMG-20261022-WA0054.jpg}}\end{tcolorbox}
\begin{tcolorbox}[title={Adım adım çözüm},colback=green!3!white,colframe=green]\begin{enumerate}
\item İlkinde \(y=0\) yolu 1, \(x=0\) yolu -1 verir. İkinci ifadede pozitif x ekseni 1, negatif x ekseni -1 verir. İkisinin de limiti yoktur.
\item Üçüncüde y sıfır olamaz; \(y=x\) yolunda 2, \(y=2x\) yolunda \(5/4\) çıkar. Bu limit de yoktur. Tanımsız yolları test yolu olarak seçme.
\end{enumerate}\end{tcolorbox}
\Needspace{10\baselineskip}\begin{tcolorbox}[title={Örnek 4.13 | Dersteki parçalı fonksiyonun sürekliliği},colframe=accent]\(f(x,y)=\sin(xy^2)/(x^2+y^2)\) orijin dışında, \(f(0,0)=0\) olsun. Her yerde sürekli olduğunu gösteriniz.
\par\vspace{1mm}{\scriptsize\sffamily Ders fotoğrafından | \nolinkurl{IMG-20261022-WA0054.jpg}}\end{tcolorbox}
\begin{tcolorbox}[title={Adım adım çözüm},colback=green!3!white,colframe=green]\begin{enumerate}
\item Orijin dışında payda sıfır değildir, dolayısıyla süreklidir. Orijinde \(|\sin u|\le|u|\) kullanılır.
\item \(|f(x,y)|\le |x|y^2/(x^2+y^2)\le|x|\le\sqrt{x^2+y^2}\to0\). Böylece limit 0 olur ve atanmış değere eşittir.
\item Kutupsal yazım da sınırı \(r|\cos\theta|\sin^2\theta\le r\) yapar. Bu tüm açılar için geçerli olduğu için kanıttır.
\end{enumerate}\end{tcolorbox}
\Needspace{10\baselineskip}\begin{tcolorbox}[title={Örnek 4.14 | Bütün doğruların aynı sonucu vermesi yetmez},colframe=accent]\(f(x,y)=x^2y/(x^4+y^2)\) orijinde limite sahip midir?
\par\vspace{1mm}{\scriptsize\sffamily Pekiştirme}\end{tcolorbox}
\begin{tcolorbox}[title={Adım adım çözüm},colback=green!3!white,colframe=green]\begin{enumerate}
\item \(y=mx\) ve sabit \(m\ne0\) için \(f=mx/(x^2+m^2)\to0\). \(m=0\) ve dikey eksen de 0 verir.
\item Fakat \(y=x^2\) boyunca \(f=1/2\). Bu nedenle limit yoktur. Bütün doğruların kontrol edilmesi bile genel limit için yeterli değildir.
\end{enumerate}\end{tcolorbox}
\clearpage\section{Kısmi türev, teğet düzlem ve lineer yaklaşım}
\Needspace{10\baselineskip}\begin{tcolorbox}[title={Örnek 5.1 | Dersteki temel kısmi türev},colframe=accent]\(f(x,y)=x^2y^3+3x\) için birinci kısmi türevleri bulunuz.
\par\vspace{1mm}{\scriptsize\sffamily Ders fotoğrafından | \nolinkurl{IMG-20261022-WA0055.jpg}}\end{tcolorbox}
\begin{tcolorbox}[title={Adım adım çözüm},colback=green!3!white,colframe=green]\begin{enumerate}
\item x yönünde y sabittir: \(f_x=2xy^3+3\). y yönünde x sabittir: \(f_y=3x^2y^2\).
\item Örneğin \((1,2)\) noktasında \(f_x=19\), \(f_y=12\). Bunlar yüzeyin koordinat kesitlerindeki eğimleridir.
\end{enumerate}\end{tcolorbox}
\Needspace{10\baselineskip}\begin{tcolorbox}[title={Örnek 5.2 | Limit tanımıyla kısmi türev},colframe=accent]\(f(x,y)=x^2y\) için \(f_x(a,b)\) değerini tanımdan bulunuz.
\par\vspace{1mm}{\scriptsize\sffamily Pekiştirme}\end{tcolorbox}
\begin{tcolorbox}[title={Adım adım çözüm},colback=green!3!white,colframe=green]\begin{enumerate}
\item \([f(a+h,b)-f(a,b)]/h=[(a+h)^2b-a^2b]/h=2ab+bh\).
\item h sıfıra giderken \(f_x(a,b)=2ab\) elde edilir. Bu, diğer değişkeni neden sabit tuttuğumuzu doğrudan gösterir.
\end{enumerate}\end{tcolorbox}
\Needspace{10\baselineskip}\begin{tcolorbox}[title={Örnek 5.3 | Dersteki rasyonel fonksiyonun ikinci türevleri},colframe=accent]\(f(x,y)=xy/(x^2+y^2)\), \((x,y)\ne(0,0)\) için ikinci kısmi türevleri bulunuz.
\par\vspace{1mm}{\scriptsize\sffamily Ders fotoğrafından | \nolinkurl{IMG-20261022-WA0055.jpg}}\end{tcolorbox}
\begin{tcolorbox}[title={Adım adım çözüm},colback=green!3!white,colframe=green]\begin{enumerate}
\item \(S=x^2+y^2\) kısaltmasıyla bölüm kuralından \(f_x=y(y^2-x^2)/S^2\), \(f_y=x(x^2-y^2)/S^2\).
\item \(f_{xx}\) için \(-2xy/S^2-4xy(y^2-x^2)/S^3\) ortak paydada toplanır: \(f_{xx}=2xy(x^2-3y^2)/S^3\).
\item Simetriyle \(f_{yy}=2xy(y^2-3x^2)/S^3\). Karma türev \(f_{xy}=(-x^4+6x^2y^2-y^4)/S^3\).
\item Orijin dışında türevler süreklidir; \(f_{yx}=f_{xy}\). Orijinde başlangıç fonksiyonu tanımlı olmadığı için bu formüller orijine uygulanmaz.
\end{enumerate}\end{tcolorbox}
\Needspace{10\baselineskip}\begin{tcolorbox}[title={Örnek 5.4 | Kısmi türevler var, süreklilik yok},colframe=accent]\(f(x,y)=xy/(x^2+y^2)\) orijin dışında, \(f(0,0)=0\) olsun. Orijinde kısmi türevleri ve türevlenebilirliği inceleyiniz.
\par\vspace{1mm}{\scriptsize\sffamily Pekiştirme}\end{tcolorbox}
\begin{tcolorbox}[title={Adım adım çözüm},colback=green!3!white,colframe=green]\begin{enumerate}
\item Eksenler boyunca fonksiyon sıfırdır. Bu yüzden \(f_x(0,0)=f_y(0,0)=0\).
\item \(y=x\) boyunca değer \(1/2\), dolayısıyla fonksiyon orijinde sürekli değildir. Sürekli olmayan fonksiyon türevlenebilir olamaz.
\end{enumerate}\end{tcolorbox}
\Needspace{10\baselineskip}\begin{tcolorbox}[title={Örnek 5.5 | Teğet düzlem ve lineer model},colframe=accent]\(f(x,y)=x^2+y^2\) için \((1,2,5)\) noktasındaki teğet düzlemi ve normal doğruyu bulunuz.
\par\vspace{1mm}{\scriptsize\sffamily Pekiştirme}\end{tcolorbox}
\begin{tcolorbox}[title={Adım adım çözüm},colback=green!3!white,colframe=green]\begin{enumerate}
\item \(f_x=2x,f_y=2y\). Noktadaki eğimler 2 ve 4. Düzlem \(z-5=2(x-1)+4(y-2)\), yani \(z=2x+4y-5\).
\item Normal doğru \((x,y,z)=(1,2,5)+t(-2,-4,1)\). Normalin işaretini ters almak aynı doğruyu verir.
\end{enumerate}\end{tcolorbox}
\Needspace{10\baselineskip}\begin{tcolorbox}[title={Örnek 5.6 | Dersteki üstel fonksiyonun iki lineerleştirmesi},colframe=accent]\(f(x,y)=e^{2y-x}\) için \(P=(0,0)\) ve \(Q=(1,2)\) merkezlerindeki lineerleştirmeleri bulunuz.
\par\vspace{1mm}{\scriptsize\sffamily Ders fotoğrafından | \nolinkurl{IMG-20261105-WA0038.jpg}}\end{tcolorbox}
\begin{tcolorbox}[title={Adım adım çözüm},colback=green!3!white,colframe=green]\begin{enumerate}
\item \(f_x=-e^{2y-x}\), \(f_y=2e^{2y-x}\). P merkezinde \(L_P=1-x+2y\).
\item Q merkezinde \(f=e^3,f_x=-e^3,f_y=2e^3\). Bu yüzden \(L_Q=e^3-e^3(x-1)+2e^3(y-2)=e^3(2y-x-2)\).
\item Her model kendi merkezinin yakınında kullanılır. Örneğin \(f(0.01,0.02)=e^{0.03}\approx1.03\); bu hedef P merkezine yakındır.
\end{enumerate}\end{tcolorbox}
\Needspace{10\baselineskip}\begin{tcolorbox}[title={Örnek 5.7 | Karekökün yaklaşık değeri},colframe=accent]\(\sqrt{(3.01)^2+(3.98)^2}\) ifadesini diferansiyel veya lineer modelle yaklaşık hesaplayınız.
\par\vspace{1mm}{\scriptsize\sffamily Pekiştirme}\end{tcolorbox}
\begin{tcolorbox}[title={Adım adım çözüm},colback=green!3!white,colframe=green]\begin{enumerate}
\item \(f(x,y)=\sqrt{x^2+y^2}\), merkez \((3,4)\), \(f=5\). \(f_x=3/5,f_y=4/5\).
\item \(h=0.01,k=-0.02\). Yaklaşık değer \(5+(3/5)(0.01)+(4/5)(-0.02)=4.99\).
\end{enumerate}\end{tcolorbox}
\Needspace{10\baselineskip}\begin{tcolorbox}[title={Örnek 5.8 | Defterdeki hata sınırı},colframe=accent]\(f=x^2+3xy+2y^2+4\), merkez \((2,1)\), \(|x-2|,|y-1|\le0.01\). Lineerleştirmeyi ve hata sınırını bulunuz.
\par\vspace{1mm}{\scriptsize\sffamily Ders fotoğrafından | \nolinkurl{IMG-20261022-WA0060.jpg}}\end{tcolorbox}
\begin{tcolorbox}[title={Adım adım çözüm},colback=green!3!white,colframe=green]\begin{enumerate}
\item \(f(2,1)=16\), \(f_x(2,1)=7\), \(f_y(2,1)=10\). \(L=16+7(x-2)+10(y-1)=7x+10y-8\).
\item \(f_{xx}=2,f_{xy}=3,f_{yy}=4\). M=4 seçilir; hata en fazla \(\tfrac42(0.01+0.01)^2=0.0008\).
\item Kontrol: gerçek kalan \(h^2+3hk+2k^2\) olur. Mutlak değer için daha sıkı \(0.0006\) sınırı da mümkündür; dersteki \(0.0008\) güvenli ve doğrudur.
\end{enumerate}\end{tcolorbox}
\Needspace{10\baselineskip}\begin{tcolorbox}[title={Örnek 5.9 | İkinci hata sınırı sorusu},colframe=accent]\(f=x^2-3xy+y^2\), merkez \((1,2)\), \(|x-1|,|y-2|\le0.1\). L ve güvenli hata sınırını bulunuz.
\par\vspace{1mm}{\scriptsize\sffamily Ders fotoğrafından | \nolinkurl{IMG-20261105-WA0038.jpg}}\end{tcolorbox}
\begin{tcolorbox}[title={Adım adım çözüm},colback=green!3!white,colframe=green]\begin{enumerate}
\item \(f(1,2)=-1,f_x=-4,f_y=1\) olduğundan \(L=-1-4(x-1)+(y-2)=-4x+y+1\).
\item İkinci türevlerin mutlak değerleri 2,3,2; M=3. Hata \(\le(3/2)(0.1+0.1)^2=0.06\).
\end{enumerate}\end{tcolorbox}
\Needspace{10\baselineskip}\begin{tcolorbox}[title={Örnek 5.10 | Küçük bölgede kosinüs hatası},colframe=accent]\(f(x,y)=1+x+\cos y\), merkez \((0,0)\), \(|x|,|y|\le0.2\). L ve hata sınırını bulunuz.
\par\vspace{1mm}{\scriptsize\sffamily Ders fotoğrafından | \nolinkurl{IMG-20261105-WA0040.jpg}}\end{tcolorbox}
\begin{tcolorbox}[title={Adım adım çözüm},colback=green!3!white,colframe=green]\begin{enumerate}
\item \(f(0,0)=2,f_x=1,f_y(0,0)=0\). Böylece \(L=2+x\).
\item \(f_{xx}=f_{xy}=0,|f_{yy}|=|\cos y|\le1\). Genel sınır 0.08 verir; yalnız y yönündeki gerçek ikinci terim kullanılırsa \(|f-L|=|\cos y-1|\le y^2/2\le0.02\).
\end{enumerate}\end{tcolorbox}
\Needspace{10\baselineskip}\begin{tcolorbox}[title={Örnek 5.11 | Tahtadaki silindir hacmi değişimi},colframe=accent]Yarıçapı 3 cm, yüksekliği 10 cm olan silindirde yarıçap 0.01 cm artsın, yükseklik 0.1 cm azalsın. Hacim değişimini yaklaşık bulunuz.
\par\vspace{1mm}{\scriptsize\sffamily Ders fotoğrafından | \nolinkurl{IMG-20261108-WA0066.jpg}}\end{tcolorbox}
\begin{tcolorbox}[title={Adım adım çözüm},colback=green!3!white,colframe=green]\begin{enumerate}
\item \(V=\pi r^2h\); \(dV=2\pi rh\,dr+\pi r^2\,dh\). \(dr=0.01,dh=-0.1\).
\item \(dV=60\pi(0.01)+9\pi(-0.1)=-0.3\pi\) santimetre küp. Başlangıç hacmi \(90\pi\), yeni hacim yaklaşık \(89.7\pi\).
\item Tam değer \(\pi(3.01)^2(9.9)=89.69499\pi\). Yaklaşık ve tam değer arasındaki fark \(0.00501\pi\); yüksek mertebeli terimlerden gelir.
\end{enumerate}\end{tcolorbox}
\Needspace{10\baselineskip}\begin{tcolorbox}[title={Örnek 5.12 | Silindir için ikinci ders sorusu},colframe=accent]Yarıçap 10 cm den 9.9 cm ye, yükseklik 30 cm den 30.3 cm ye değişiyor. Yaklaşık hacim değişimi nedir?
\par\vspace{1mm}{\scriptsize\sffamily Ders fotoğrafından | \nolinkurl{IMG-20261108-WA0083.jpg}}\end{tcolorbox}
\begin{tcolorbox}[title={Adım adım çözüm},colback=green!3!white,colframe=green]\begin{enumerate}
\item \(dr=-0.1,dh=0.3\). \(dV=2\pi(10)(30)(-0.1)+\pi(100)(0.3)=-30\pi\) santimetre küp.
\item Başlangıç hacmi \(3000\pi\), yeni hacim yaklaşık \(2970\pi\). İşaret hacmin azaldığını gösterir.
\end{enumerate}\end{tcolorbox}
\Needspace{10\baselineskip}\begin{tcolorbox}[title={Örnek 5.13 | Tahtadaki polinom diferansiyeli},colframe=accent]\(f=x^3+5xy^2+3xy+4x-3y+5\) için df ve ikinci diferansiyeli yazınız.
\par\vspace{1mm}{\scriptsize\sffamily Ders fotoğrafından | \nolinkurl{IMG-20261108-WA0082.jpg}}\end{tcolorbox}
\begin{tcolorbox}[title={Adım adım çözüm},colback=green!3!white,colframe=green]\begin{enumerate}
\item \(f_x=3x^2+5y^2+3y+4\), \(f_y=10xy+3x-3\). df bu katsayılarla dx ve dy nin toplamıdır.
\item \(f_{xx}=6x,f_{xy}=10y+3,f_{yy}=10x\). Dolayısıyla \(d^2f=6x\,dx^2+2(10y+3)dx\,dy+10x\,dy^2\).
\end{enumerate}\end{tcolorbox}
\Needspace{10\baselineskip}\begin{tcolorbox}[title={Örnek 5.14 | Üç değişkenli polinomun diferansiyelleri},colframe=accent]\(w=3x^2yz^2-3xy+2yz^2+xyz\) için birinci ve ikinci diferansiyeli bulunuz.
\par\vspace{1mm}{\scriptsize\sffamily Ders fotoğrafından | \nolinkurl{IMG-20261108-WA0082.jpg}}\end{tcolorbox}
\begin{tcolorbox}[title={Adım adım çözüm},colback=green!3!white,colframe=green]\begin{enumerate}
\item \(w_x=6xyz^2-3y+yz\), \(w_y=3x^2z^2-3x+2z^2+xz\), \(w_z=6x^2yz+4yz+xy\). \(dw=w_xdx+w_ydy+w_zdz\).
\item \(w_{xx}=6yz^2,w_{yy}=0,w_{zz}=6x^2y+4y\). Karma türevler \(w_{xy}=6xz^2-3+z\), \(w_{xz}=12xyz+y\), \(w_{yz}=6x^2z+4z+x\).
\item İkinci diferansiyel bu altı katsayı yukarıdaki genel formüle yerleştirilerek elde edilir. Her karma terimin katsayısı 2 ile çarpılır.
\end{enumerate}\end{tcolorbox}
\clearpage\section{Zincir kuralı, yönlü türev ve gradyan}
\Needspace{10\baselineskip}\begin{tcolorbox}[title={Örnek 6.1 | Tahtadaki trigonometrik zincir örneği},colframe=accent]\(z=x^2+y^2\), \(x=\cos t+\sin t\), \(y=\cos t-\sin t\). dz/dt nedir?
\par\vspace{1mm}{\scriptsize\sffamily Ders fotoğrafından | \nolinkurl{IMG-20261105-WA0041.jpg}}\end{tcolorbox}
\begin{tcolorbox}[title={Adım adım çözüm},colback=green!3!white,colframe=green]\begin{enumerate}
\item \(z_x=2x,z_y=2y\); \(x\prime=-\sin t+\cos t\), \(y\prime=-\sin t-\cos t\). Zincir toplamı 0 olur.
\item Doğrudan kontrol: \((\cos t+\sin t)^2+(\cos t-\sin t)^2=2\). Sabit fonksiyonun türevi 0.
\end{enumerate}\end{tcolorbox}
\Needspace{10\baselineskip}\begin{tcolorbox}[title={Örnek 6.2 | Çarpanlı fonksiyonda zincir kuralı},colframe=accent]\(z=xy\), \(x=\cos t,y=\sin t\). dz/dt yi bulunuz.
\par\vspace{1mm}{\scriptsize\sffamily Ders fotoğrafından | \nolinkurl{IMG-20261105-WA0041.jpg}}\end{tcolorbox}
\begin{tcolorbox}[title={Adım adım çözüm},colback=green!3!white,colframe=green]\begin{enumerate}
\item \(z_x=y,z_y=x\). Böylece \(z\prime=\sin t(-\sin t)+\cos t\cos t=\cos^2t-\sin^2t\).
\item Doğrudan \(z=\tfrac12\sin2t\) türevi de \(\cos2t\), aynı sonuçtur.
\end{enumerate}\end{tcolorbox}
\Needspace{10\baselineskip}\begin{tcolorbox}[title={Örnek 6.3 | Tahtadaki logaritma örneği},colframe=accent]\(z=\ln(x^2+y^2)\), \(x=e^u\cos v,y=e^u\sin v\). \(z_u,z_v\) nedir?
\par\vspace{1mm}{\scriptsize\sffamily Ders fotoğrafından | \nolinkurl{IMG-20261108-WA0074.jpg}}\end{tcolorbox}
\begin{tcolorbox}[title={Adım adım çözüm},colback=green!3!white,colframe=green]\begin{enumerate}
\item \(f_x=2x/(x^2+y^2)\), \(f_y=2y/(x^2+y^2)\). \(x_u=x,y_u=y,x_v=-y,y_v=x\).
\item \(z_u=2(x^2+y^2)/(x^2+y^2)=2\); \(z_v=(-2xy+2xy)/(x^2+y^2)=0\).
\item Kontrol: \(x^2+y^2=e^{2u}\), dolayısıyla \(z=2u\).
\end{enumerate}\end{tcolorbox}
\Needspace{10\baselineskip}\begin{tcolorbox}[title={Örnek 6.4 | Üç ara değişkenli ders örneği},colframe=accent]\(S=\ln(x^2+y^2+z^2)\), \(x=u+v,y=u-v,z=2\sqrt{uv}\). Pozitif u ve v için \(S_u,S_v\) bulunuz.
\par\vspace{1mm}{\scriptsize\sffamily Ders fotoğrafından | \nolinkurl{IMG-20261108-WA0077.jpg}}\end{tcolorbox}
\begin{tcolorbox}[title={Adım adım çözüm},colback=green!3!white,colframe=green]\begin{enumerate}
\item \(x_u=x_v=1,y_u=1,y_v=-1,z_u=v/\sqrt{uv},z_v=u/\sqrt{uv}\).
\item Payda \(x^2+y^2+z^2=2(u+v)^2\). Zincir toplamları \(S_u=S_v=2/(u+v)\).
\item Doğrudan \(S=\ln2+2\ln(u+v)\) aynı sonucu verir. \(u,v>0\) seçimi kökün içte türevlenmesini ve logaritmanın tanımını güvenceye alır.
\end{enumerate}\end{tcolorbox}
\Needspace{10\baselineskip}\begin{tcolorbox}[title={Örnek 6.5 | İki bağımsız yeni değişkenle karmaşık örnek},colframe=accent]\(z=e^{x/y}\), \(x=\sin u,y=u^2v\), \(u,v\ne0\). \(z_u,z_v\) yi bulunuz.
\par\vspace{1mm}{\scriptsize\sffamily Ders fotoğrafından | \nolinkurl{IMG-20261108-WA0079.jpg}}\end{tcolorbox}
\begin{tcolorbox}[title={Adım adım çözüm},colback=green!3!white,colframe=green]\begin{enumerate}
\item \(f_x=e^{x/y}/y\), \(f_y=-xe^{x/y}/y^2\). \(x_u=\cos u,x_v=0,y_u=2uv,y_v=u^2\).
\item \(z_u=e^{\sin u/(u^2v)}[\cos u/(u^2v)-2\sin u/(u^3v)]\).
\item \(z_v=-e^{\sin u/(u^2v)}\sin u/(u^2v^2)\). İç ifadeyi önce sadeleştirip türev almak da kontrol yöntemidir.
\end{enumerate}\end{tcolorbox}
\Needspace{10\baselineskip}\begin{tcolorbox}[title={Örnek 6.6 | Tahtadaki kapalı türev sorusu},colframe=accent]\(2x^4z^3+3xy^2+4xz+3y+5x-23=0\) denkleminin tanımladığı z(x,y) için \((1,2)\) noktasında \(z_x,z_y\) nedir?
\par\vspace{1mm}{\scriptsize\sffamily Ders fotoğrafından | \nolinkurl{IMG-20261108-WA0081.jpg}}\end{tcolorbox}
\begin{tcolorbox}[title={Adım adım çözüm},colback=green!3!white,colframe=green]\begin{enumerate}
\item Önce z bulunur: \(2z^3+4z=0\). Gerçek çözüm \(z=0\); inceleme noktası \((1,2,0)\).
\item \(F_x=8x^3z^3+3y^2+4z+5\), \(F_y=6xy+3\), \(F_z=6x^4z^2+4x\).
\item Noktada \(F_x=17,F_y=15,F_z=4\ne0\). Böylece \(z_x=-17/4,z_y=-15/4\). Yerel z(x,y) yorumu geçerlidir.
\end{enumerate}\end{tcolorbox}
\Needspace{10\baselineskip}\begin{tcolorbox}[title={Örnek 6.7 | Kapalı denklemin teğet düzlemi},colframe=accent]\(x^2+y^2+z^2=9\) küresinin \((1,2,2)\) noktasındaki teğet düzlemi nedir?
\par\vspace{1mm}{\scriptsize\sffamily Pekiştirme}\end{tcolorbox}
\begin{tcolorbox}[title={Adım adım çözüm},colback=green!3!white,colframe=green]\begin{enumerate}
\item \(F=x^2+y^2+z^2-9\), gradyanı \((2x,2y,2z)\). Noktada normal \((2,4,4)\).
\item \(2(x-1)+4(y-2)+4(z-2)=0\), yani \(x+2y+2z=9\). Bu normal merkezden noktaya giden yarıçapla paraleldir.
\end{enumerate}\end{tcolorbox}
\Needspace{10\baselineskip}\begin{tcolorbox}[title={Örnek 6.8 | Dersteki diferansiyel denklem dönüşümü},colframe=accent]\(x=e^t\) dönüşümü altında \(x^2y_{xx}+xy_x-y=0\) denklemini t değişkenine dönüştürünüz.
\par\vspace{1mm}{\scriptsize\sffamily Ders fotoğrafından | \nolinkurl{IMG-20261105-WA0042.jpg}}\end{tcolorbox}
\begin{tcolorbox}[title={Adım adım çözüm},colback=green!3!white,colframe=green]\begin{enumerate}
\item \(dt/dx=1/x\). Dolayısıyla \(y_x=y_t/x\).
\item Yeniden x türevi: \(y_{xx}=(y_{tt}-y_t)/x^2\). Burada hem \(y_t\) hem \(1/x\) türevlenir.
\item Denklem \((y_{tt}-y_t)+y_t-y=0\), yani \(y_{tt}-y=0\) olur. Dönüşüm \(x>0\) için geçerlidir.
\end{enumerate}\end{tcolorbox}
\Needspace{10\baselineskip}\begin{tcolorbox}[title={Örnek 6.9 | Birim olmayan yönden yönlü türev},colframe=accent]\(f(x,y)=x^2y+3y\), \(P=(1,2)\), yön \(\mathbf v=(3,4)\). Yönlü türev nedir?
\par\vspace{1mm}{\scriptsize\sffamily Pekiştirme}\end{tcolorbox}
\begin{tcolorbox}[title={Adım adım çözüm},colback=green!3!white,colframe=green]\begin{enumerate}
\item \(f_x=2xy,f_y=x^2+3\). \(\nabla f(1,2)=(4,4)\).
\item Yönün normu 5; \(\mathbf u=(3/5,4/5)\). \(D_{\mathbf u}f=4(3/5)+4(4/5)=28/5\).
\item Birimleştirmeden alınan 28 sonucu, parametre hızını 5 katına çıkarır ve bu sorunun istediği birim uzunluk başına türev değildir.
\end{enumerate}\end{tcolorbox}
\Needspace{10\baselineskip}\begin{tcolorbox}[title={Örnek 6.10 | P den Q ya doğru değişim},colframe=accent]\(F(x,y,z)=x^2+y^2+z^2\), \(P=(1,0,2)\), \(Q=(3,1,4)\). P de Q yönündeki türev nedir?
\par\vspace{1mm}{\scriptsize\sffamily Pekiştirme}\end{tcolorbox}
\begin{tcolorbox}[title={Adım adım çözüm},colback=green!3!white,colframe=green]\begin{enumerate}
\item \(Q-P=(2,1,2)\), norm 3; birim yön \((2,1,2)/3\). \(\nabla F(P)=(2,0,4)\).
\item Skaler çarpım \((4+0+8)/3=4\). Sonucun pozitif olması bu yönde ilk anda alan değerinin arttığını gösterir.
\end{enumerate}\end{tcolorbox}
\Needspace{10\baselineskip}\begin{tcolorbox}[title={Örnek 6.11 | Gradyanla yön seçmek},colframe=accent]\(f=x^2+2y^2\) için P=(1,1) noktasında en hızlı artış yönü, hızı ve sıfır değişim yönü nedir?
\par\vspace{1mm}{\scriptsize\sffamily Pekiştirme}\end{tcolorbox}
\begin{tcolorbox}[title={Adım adım çözüm},colback=green!3!white,colframe=green]\begin{enumerate}
\item \(\nabla f=(2x,4y)\); P de \((2,4)\), norm \(2\sqrt5\). En hızlı artış birim yönü \((1,2)/\sqrt5\), hızı \(2\sqrt5\).
\item En hızlı azalış ters yönde ve türev \(-2\sqrt5\). Sıfır türev yönlerinden biri \((-2,1)/\sqrt5\), çünkü \((2,4)\cdot(-2,1)=0\).
\end{enumerate}\end{tcolorbox}
\Needspace{10\baselineskip}\begin{tcolorbox}[title={Örnek 6.12 | Seviye eğrisinin teğet doğrusu},colframe=accent]\(x^2+2y^2=3\) eğrisinin (1,1) noktasındaki teğeti nedir?
\par\vspace{1mm}{\scriptsize\sffamily Pekiştirme}\end{tcolorbox}
\begin{tcolorbox}[title={Adım adım çözüm},colback=green!3!white,colframe=green]\begin{enumerate}
\item Gradyan \((2,4)\) teğete normaldir. \(2(x-1)+4(y-1)=0\), yani \(x+2y=3\).
\item Bu doğrultu önceki örnekteki sıfır yönlü türev doğrultusuyla aynıdır. Teğet doğrusu eğriye yalnızca yerel olarak yaklaşır.
\end{enumerate}\end{tcolorbox}
\Needspace{10\baselineskip}\begin{tcolorbox}[title={Örnek 6.13 | Birleşik pekiştirme},colframe=accent]\(f(x,y)=e^{x-y}\), P=(0,0). P de teğet düzlem, \((1,1)\) yönünde türev ve en hızlı artış yönünü bulunuz.
\par\vspace{1mm}{\scriptsize\sffamily Pekiştirme}\end{tcolorbox}
\begin{tcolorbox}[title={Adım adım çözüm},colback=green!3!white,colframe=green]\begin{enumerate}
\item \(f(P)=1,\nabla f(P)=(1,-1)\). Grafik üzerindeki teğet düzlem \(z=1+x-y\).
\item Birim yön \((1,1)/\sqrt2\); yönlü türev 0. En hızlı artış yönü \((1,-1)/\sqrt2\), artış hızı \(\sqrt2\).
\item \(x=y\) doğrusu üzerinde fonksiyon gerçekten 1 dir. Bu örnekte sıfır türev sabit seviye doğrultusuna karşılık gelir; her fonksiyonda böyle küresel bir sonuç çıkarılmaz.
\end{enumerate}\end{tcolorbox}

\clearpage
\section*{Kaynaklar ve kapsam}
\addcontentsline{toc}{section}{Kaynaklar ve kapsam}
Bu metin 156 ders/defter fotoğrafı ve altı başlıklı müfredata göre özgün olarak yazıldı. Fotoğraflar dersin sıra ve örneklerini belirler. Ek sorular pekiştirme olarak etiketlidir; çıkmış sınav sorusu iddiası yoktur. Soluk semboller tahminle aktarılmadı. Gradyan ve yönlü türev, müfredat gereği tamamlanan anlatımdır. Fotoğraf bazında inceleme dizini web sitesinde ve analysis/FOTOGRAF\_INCELEMESI.md dosyasında bulunur.
\textbf{Birincil ders malzemesi}. Kullanıcının sağladığı 156 JPG ders/defter fotoğrafı ve mufredat.png. Dosya tarihleri ders tarihi olarak varsayılmadı. Konu-fotoğraf eşleştirmesi fotoğraf dizinindedir.\par
\textbf{Erwin Kreyszig · Advanced Engineering Mathematics}. Dersin belirtilen kitabı. Wiley 10. baskı bölüm diziniyle vektör diferansiyel hesabı başlığı eşleştirildi; tam kitap dosyası sağlanmadı. Kitaptan birebir metin veya soru aktarımı yapılmadı.
{\footnotesize\url{https://bcs.wiley.com/he-bcs/Books?action=contents&bcsId=6140&itemId=0470458364}}\par
\textbf{KOCW · Kreyszig temelli Bölüm 9 ders notları}. Sungkyunkwan / Kim Munhyeon, vektörler ve vektör fonksiyonları. Açık üniversite notu kavram ve formül kontrolünde kullanıldı. Buradaki Türkçe açıklamalar ve çözümler özgün olarak yazıldı.
{\footnotesize\url{https://contents.kocw.net/KOCW/document/2015/sungkyunkwan/kimmunhyeon3/8.pdf}}\par
\textbf{Zill ve Wright · Calculus / Matematik}. Kullanıcının verdiği bibliyografya: İ. N. Cangül çevirisi, Nobel, 2013. Tam metne erişilmedi; baskıya özgü sayfa ve soru numarası atfedilmedi.\par
\textbf{Mustafa Balcı · Genel Matematik II}. Kullanıcının verdiği bibliyografya: Sürat, 2013. Tam metin sağlanmadı; bu kitaptan doğrulanmış alıntı veya sayfa atfı yapılmadı.\par
\textbf{MIT OCW · 18.02 ders notları}. Denis Auroux, 2007. Ders düzeni ve kavram denetimi için açık ders dizini. Müfredat dışındaki çok katlı integraller ve integral teoremleri bu çalışmaya eklenmedi.
{\footnotesize\url{https://ocw.mit.edu/courses/18-02-multivariable-calculus-fall-2007/pages/lecture-notes/}}\par
\textbf{MIT · Teğet düzlem ve gradyan}. Gradyanın teğet düzlem ve yönlü türevle ilişkisi için açık teknik kontrol kaynağı. Konuya eklenen örnekler özgün pekiştirme sorularıdır.
{\footnotesize\url{https://www.ocw.mit.edu/ans7870/18/18.013a/textbook/HTML/chapter06/section05.html}}\par
\end{document}
